\subsection{Fourier iteration with no prescribed midpoint}
\label{uf:transfer}

We use the regularized component \(\sigma_i\), the conditional energy
estimate \eqref{fprof:energy}, and the deletion estimates
\eqref{fprof:sourcebad}--\eqref{fprof:pinbad}. The chain required in this
subsection has the following properties:
\begin{gather}\label{uf:chain}
 n_0>n_1>\cdots>n_K=0,\qquad 2n_j-n_{j+1}\leq N,\qquad
 n_0=(1-\zeta)N+O(T),\\
 K\leq K_0,\qquad
 \sum_{j<K}c_{i,j}\leq\mathcal C N+C K\Gamma N+C_{T,K},\qquad
 c_{i,j}=g_i(n_{j+1})-\min_{[n_{j+1},n_j]}g_i.\nonumber
\end{gather}
Here \(0<\zeta<1/8\), \(K_0\) is fixed independently of \(R=2^N\),
and the depths lie on the \(T\)-grid. Here \(\mathcal C\) is the cost bound in
Proposition~\ref{human:transfer}. No intermediate depth is prescribed.
We retain \(N\in4T\mathbb N\), so the standard angular scale
used below is dyadic.

Put
$$
 a_j=2^{-n_j},\qquad r_j=Ra_j,\qquad
 \delta_j=r_j^{-1},\qquad
 \rho_j=\frac{a_j}{a_{j+1}},\qquad
 \delta_*=R^{-1/2}.
$$
The two elementary inequalities
\begin{equation}\label{uf:angles}
 \delta_j\leq\rho_j,\qquad \delta_*\leq\rho_j
 \qquad(0\leq j<K)
\end{equation}
have different roles. The first is equivalent to
\(r_{j+1}\leq r_j^2\). For the second, admissibility and
\(n_{j+1}\geq0\) give
$$
 n_j-n_{j+1}\leq\min\{n_j,N-n_j\}\leq N/2.
$$
The last edge also gives \(n_{K-1}\leq N/2\), hence
\(\delta_{K-1}\leq\delta_*\). The second inequality in
\eqref{uf:angles} permits the standard source wave packets to be retained
even when the auxiliary Fourier caps become narrower than \(\delta_*\).

\subsubsection*{A fixed source packet decomposition and a separate label tree}

Fix a smooth annular multiplier at the scale \(R\). First partition it
into the usual smooth standard-cap multipliers \(\psi_{R,S}\), with
angular width \(\delta_*\). The label \(S\) has a nominal dyadic angular
arc of that length; its multiplier is supported in a fixed enlargement
of the arc, and these supports have bounded overlap. These standard
multipliers are fixed before any finer partition is made.

Use their usual source wave packets, with spatial tubes of width
\(R^{-1/2+h}\) and fixed bounded length. Choose the longitudinal span
to cover the common normalized source and pin region, with its cutoff
outside that region. Thus remoteness there is transverse separation,
not an artifact of a tube endpoint. Write
$$
 M_T\mu=\chi\,\eta_T\,(\psi_{R,S(T)}\widehat\mu)^\vee.
$$
Here \(\chi\) is a fixed smooth compact source cutoff, equal to one on
a neighborhood of \(\supp\mu\), and with support separated from the pin
set. For each standard label the smooth spatial functions \(\eta_T\)
form a partition of unity on \(\supp\chi\). Thus summing all standard
packets gives precisely \(P_R\mu\), with the source cutoff convention
used above. All packet localization estimates are applied only to these
smooth standard multipliers and these standard spatial tubes.

We now define a separate angular tree for Fourier estimates. For each
level with \(\delta_j\geq\delta_*\), its labels group whole standard
labels according to their nominal dyadic ancestors of width \(\delta_j\).
Their multipliers are sums of the corresponding \(\psi_{R,S}\).
For \(\delta_j<\delta_*\), a level label is a pair \((S,I)\), where
\(I\) is a dyadic angular interval of length \(\delta_j\) meeting the
support of \(\psi_{R,S}\); its multiplier is
$$
 \psi_{R,S}(\xi)\,\mathbf 1_I(\xi/|\xi|).
$$
The intervals partition the entire angular support of the standard
multiplier, including the part outside its nominal arc. Consequently
each such interval is within angular distance \(O(\delta_*)\) of the
standard label's direction. Use half-open angular intervals to specify
the partition at endpoints. The endpoint choices do not affect any
circular integral.

After the standard scale, parents keep the same label \(S\) and group
the finer dyadic intervals. Before it, parents group whole standard
labels. If an edge of the spatial chain crosses the standard angular
scale, insert that scale in the angular tree only. This creates no
additional spatial level and no additional inflation factor.
The notation \(\theta\prec\alpha\) below means descent in this label
tree. It is a relation between labels; overlapping smooth standard
supports cause no ambiguity in this relation.

Every level-\(j\) multiplier has angular support of width
\(O(\delta_j)\), and the supports have bounded overlap at each level.
For fine labels this follows because at each direction only boundedly
many standard multipliers are nonzero, and the dyadic intervals partition
that direction separately for each such multiplier. For coarse labels
it follows by grouping adjacent standard labels. No differentiability
of a fine angular indicator will be used: these indicators partition
spectral measures on circles, not source packets.

\subsubsection*{Coarsened marks and their bad-pair estimate}

For a level-\(j\) Fourier label \(\theta\), define its effective label
to be \(\theta\) itself if \(\delta_j\geq\delta_*\), and its standard
parent \(S\) otherwise. Its angular width is
$$
 \overline\delta_j=\max\{\delta_j,\delta_*\}\leq\rho_j.
$$
For each active \(a_j\)-cube \(Q_j\subset Q_{j+1}\), mark an effective
label bad when the enlarged \(a_j\)-by-\(a_{j+1}\) test tube, centered
at \(Q_j\) and directed along that effective label, has mass greater than
\(H_{i,j}\rho_j\) under \((\sigma_i)_{Q_{j+1}^+}\).
All Fourier labels with the same effective label receive the same mark.
Descendant spatial cubes inherit the mark from their unique level-\(j\)
ancestor. Use the thresholds \eqref{fprof:thresholds}.

For definiteness, the averaging cubes and test dilations may be chosen as
$$
 Q_j^+=R^{(2j+2)h}Q_j,
 \qquad\text{test-tube dilation }R^{(4K_0+20)h}.
$$
For large \(R\), child averaging cubes and their further \(R^h\)
localizations fit inside the corresponding enlarged parent. Additional
fixed geometric factors are absorbed into these dilations or their
constants. The same enlarged parent is used in the marking condition,
the truncated conditional energy, and the Fourier iteration.

These marks are unions of whole standard source-cap labels. At coarse
levels this is true by their definition as groups of standard labels;
at fine levels all descendants of \(S\) share the mark of \(S\).
For an initial pin cube \(Q_0\), remove a standard source packet when
its enlarged tube \(2T\) meets \(Q_0\) and its standard label has any
inherited bad mark. Call the retained source function \(G_{R,i,Q_0}\).

The deletion estimates already proved apply to these marks. Indeed,
directions within \(O(\overline\delta_j)\) of an effective direction
are within \(O(\rho_j)\), so the same \(\rho_j\)-net of directions
and the same containing tube grid cover the marked tubes. Their widths
and conditional mass normalization are exactly those used in
\eqref{fprof:energy}--\eqref{fprof:pinbad}, up to the permitted
enlargements. Source-heavy directions remain only an auxiliary device
for estimating the bad mass; they are not additional Fourier marks.

The standard packet estimates give
\(\|M_T\mu\|_1\lesssim\mu(CT)+\operatorname{RapDec}(R)\),
and rapidly small pinned first norm outside an enlarged \(T\).
A separated source-pin pair lies in only \(R^{O(h)}\) enlarged
standard packet tubes: their directions are within
\(O(R^{-1/2+h})\) of the connecting line, and the standard directions
have spacing comparable to \(R^{-1/2}\). Spatial tube multiplicity
for a fixed direction is bounded. When a removed packet contributes to
such a pair, its standard label and its effective bad label differ in
direction by \(O(\overline\delta_j)\); its packet uncertainty is
\(O(R^h\delta_*)\). By \eqref{uf:angles}, both fit inside the enlarged
conditional test tube based on that pair.

It follows, with \(\mathrm{Bad}_{R,i}\) the union of bad pairs estimated
by the deletion argument, that
\begin{equation}\label{fprof:badL1}
 \int\|(d_y)_*(P_R\mu-G_{R,i,Q_0(y)})\|_1\,d\sigma_i(y)
 \leq C R^{C_{K_0}h}(\mu\times\sigma_i)(\mathrm{Bad}_{R,i})
       +\operatorname{RapDec}(R).
\end{equation}
Only standard packets were summed here. There is no multiplicity equal
to the number of auxiliary fine Fourier labels in a standard cap.
Taking \(D_1,D_2\) sufficiently large in terms of \(q,K_0\), and then
\(\Gamma\) sufficiently small compared with \(h\), makes the right
side at most \(CR^{-c}\) for some \(c>0\), uniformly in the component.

For later reference, put \(G_{R,y}=G_{R,i,Q_0(y)}\) on the component
\(X_i\), and zero on the discarded pin remainder. Multiplying by the
component weights and summing proves
\begin{equation}\label{uf:jointL1}
 \int\|(d_y)_*(P_R\mu-G_{R,y})\|_1\,d\nu(y)\leq CR^{-c'}
 \qquad(c'>0).
\end{equation}
On the discarded remainder use \(\|P_R\mu\|_1\leq C_T\) and its
mass at most \(R^{-\epsilon}\). This uses the first norm of the whole
annular multiplier, rather than a sum of absolute packet norms there.

\subsubsection*{Initial reconstruction and the exact inherited identity}

Fix a circle radius \(r\asymp_T R\), and let \(\omega_r\) be its
normalized arclength measure. Choose a nonnegative compactly supported
smooth frequency function \(\psi\) whose inverse transform is bounded
away from zero on the fixed pin region. For every terminal Fourier
label \(\ell\), with multiplier \(\psi_{R,\ell}\) specified by the
tree above, set
$$
 F_\ell=\big((\psi_{R,\ell}\widehat\mu\,d\omega_r)*\psi\big)^\vee.
$$
The terminal label is a standard label if its width equals
\(\delta_*\), and a pair \((S,I)\) otherwise. The fixed frequency
convolution makes these functions smooth, even when the fine angular
partition uses indicators. Define coarser functions by summing terminal
descendants. In particular, summing all terminal descendants of \(S\)
gives the original standard-cap function exactly.

On \(Q_0\), the circular extension of the retained standard packets,
multiplied by \(\check\psi\), agrees up to a rapid error with the sum
of the retained full standard-cap functions. To verify this statement,
sum all packets of each standard label first. A packet with
\(2T\cap Q_0=\varnothing\) has rapidly small circular extension on
\(Q_0\), by the standard integration-by-parts packet estimate used in
Liu's Lemma 4.2. Therefore retaining those remote packets from a bad
standard label changes the extension there by only a rapid error.

The fixed source cutoff does not change this conclusion. The kernel of
a smooth standard annular-cap multiplier is concentrated at transverse
scale \(R^{-1/2}\) and longitudinal scale \(R^{-1}\). Outside a fixed
neighborhood of \(\supp\mu\), that multiplier applied to \(\mu\),
together with any prescribed finite number of derivatives, is rapidly
small in integrable norms. Consequently replacing
\(\chi(\psi_{R,S}\widehat\mu)^\vee\) by
\((\psi_{R,S}\widehat\mu)^\vee\) changes its circular extension by a
rapid error, uniformly for the radii in question. All source functions
used for pinned densities still include \(\chi\), preserving their
fixed compact support and source-pin separation.

All terminal descendants of a standard label have the same survival
decision at \(Q_0\). Thus, only after this whole-standard-cap
reconstruction, partition each retained standard spectral measure into
its terminal Fourier descendants. It follows that
\begin{equation}\label{uf:initialidentity}
 \check\psi(y)
 \big(G_{R,i,Q_0}*\widehat\omega_r\big)(y)
 =F_{Q_0,\mathrm{root}}(y)+\operatorname{RapDec}(R),
 \qquad y\in Q_0,
\end{equation}
where the function on the right is the sum of surviving terminal labels.
No fine angular cutoff has been applied to an individual source packet.
In particular, no localization assertion about the spatial tails of such
an individual fine piece is needed.

More generally, for an active level-\(j\) cube and a parent angular
label \(\alpha\), define \(F_{Q_j,\alpha}\) to be the sum of terminal
descendants of \(\alpha\) with no inherited bad mark at levels
\(j,j+1,\ldots,K-1\). Level zero has a single parent label, the root.
At level \(K\), no marks remain, so \(F_{Q_K,\ell}=F_\ell\).
The level-\(j\) mark is constant on each level-\(j\) Fourier label;
future marks are inherited from the parent cube. Hence
\begin{equation}\label{fprof:identity}
 F_{Q_j,\alpha}
 =\sum_{\substack{\theta\prec\alpha\text{ at level }j\\
                    \theta\text{ good for }Q_j}}
       F_{Q_{j+1},\theta}.
\end{equation}
This is an exact partition of terminal labels. The function on the right
for a fixed parent and label is independent of the child cube. After
the standard angular scale, every future mark depends only on the
standard parent \(S\); there a fixed fine function is simply retained
or omitted as a whole. This is compatible with the exact identity.

\subsubsection*{One threshold factor per inflation step}

Let \(m_Q\) be normalized Lebesgue measure on \(Q\), and define
$$
 \mathcal E_j=
 \sum_{Q_j\ \mathrm{active}}\sigma_i(Q_j^+)
 \sum_{\alpha\text{ at level }j-1}
       \int|F_{Q_j,\alpha}|^2\,dm_{Q_j^+}.
$$
The level-zero angular sum consists only of the root. Initial
localization and \eqref{uf:initialidentity} give
$$
 \int|G_{R,i,Q_0(y)}*\widehat\omega_r(y)|^2\,d\sigma_i(y)
 \leq C_T R^{C\zeta+C_{K_0}h}\mathcal E_0
       +\operatorname{RapDec}(R).
$$
Indeed each initial Fourier function has frequencies of size \(O_T(R)\).
Its pointwise square is controlled by a normalized average on an
\(R^{-1}\)-ball with rapidly decreasing tails. Replacing that ball by
an enlarged initial cube costs at most its area ratio
\(R^{2\zeta+O(T/N)+O_{K_0}(h)}\). Multiplication by its pin cube mass
and summation give the displayed bound without any regularity assumption
on the conditional pin measure.

We prove
\begin{equation}\label{fprof:inflation}
 \mathcal E_j\leq C_T R^{C_{K_0}h}H_{i,j}\mathcal E_{j+1}
                   +\operatorname{RapDec}(R).
\end{equation}
Apply Proposition~\ref{uf:embeddingprop} to each spatial parent and
parent angular label, with \(L=R^h\), \(a=a_j\), \(b=a_{j+1}\),
\(\sigma=\sigma_i\), and \(H=H_{i,j}\). Its hypotheses are verified
in detail at the end of Appendix~\ref{uf:embedding}: admissibility controls
the curvature of the frequency arc, the coarsened marking controls its
directional mismatch, and the same enlarged parent normalizes both deletion
and embedding. Summing the proposition over parents and labels gives
\eqref{fprof:inflation}. Its error is rapidly decreasing because all relevant
counts and global norms are polynomial in \(R\), while the decay order in
the proposition is arbitrary. This step introduces one factor \(H_{i,j}\),
independently of the number of finer angular labels.

\subsubsection*{The terminal estimate and the chain cost}

Iterating \eqref{fprof:inflation}, using \eqref{uf:chain} and
\eqref{fprof:thresholds}, gives the multiplicative loss
$$
 R^{\mathcal C+C\zeta+C_{q,K_0}h+C K\Gamma}.
$$
At level \(K\), the spatial scale is one and the functions are the
terminal Fourier functions \(F_\ell\) with no remaining marks. Their
global Lebesgue norms bound the remaining unit-scale averages, up to
the allowed enlargement factors. By Plancherel and Cauchy--Schwarz,
\begin{equation}\label{fprof:circleCS}
 \sum_\ell\|F_\ell\|_2^2
 \leq C_T R^{-1}
       \int_{S^1}|\widehat\mu(r\omega)|^2\,d\omega.
\end{equation}
For clarity, the normalized mass of the circle in a fixed-radius
frequency ball is \(O_T(R^{-1})\). For each terminal label, apply
Cauchy--Schwarz to the convolution defining \(F_\ell\), using this
mass bound. Integrate in frequency and sum the squared multipliers;
their bounded overlap proves \eqref{fprof:circleCS}. This reasoning is
valid for all terminal widths \(\delta_{K-1}\leq R^{-1/2}\), including
strictly finer widths. It does not introduce their number as a factor.

We conclude
\begin{equation}\label{fprof:circleestimate}
 \int|G_{R,i,Q_0(y)}*\widehat\omega_r(y)|^2\,d\sigma_i(y)
 \leq C_T R^{-1+\mathcal C+C\zeta+C_{q,K_0}h+C K\Gamma}
       \int_{S^1}|\widehat\mu(r\omega)|^2\,d\omega
       +\operatorname{RapDec}(R).
\end{equation}
Together with \eqref{uf:jointL1}, this supplies the two shell estimates
needed in the following summation argument. Its parameter choices may
be made in the same order: first \(\zeta\), fixing \(K_0\); then the
threshold constants in terms of \(q,K_0\); then \(h\); and finally
\(T,\epsilon\), making \(\Gamma\) sufficiently small. Fixed constants
depending on the resulting annular ratio are absorbed for sufficiently
large \(R\).

In particular the proved transfer uses the cost of the unforced chain
\eqref{uf:chain}. The source packet decomposition remains at the
standard scale throughout; only the spectral functions used for
quadratic inflation are refined past that scale.
